% ECONOMICS_PROBLEM: {"id": "J31-440", "title": "Wages versus job amenities", "jel_code": "J31", "source_id": "OP-0355"}
% !TeX program = xelatex
\documentclass[11pt,a4paper]{article}
\usepackage[margin=23mm]{geometry}
\usepackage{fontspec}
\setmainfont{Latin Modern Roman}
\usepackage{amsmath,amssymb,mathtools}
\usepackage{xcolor,enumitem,booktabs,longtable,array,xurl,hyperref}
\definecolor{muted}{HTML}{53636C}
\hypersetup{colorlinks=true,urlcolor=blue,linkcolor=blue}
\setlength{\parindent}{0pt}
\setlength{\parskip}{5pt}
\newcommand{\R}{\mathbb R}
\newcommand{\E}{\mathbb E}
\newcommand{\N}{\mathbb N}
\newcommand{\ind}{\mathbf 1}
\newcommand{\Var}{\operatorname{Var}}
\newcommand{\Cov}{\operatorname{Cov}}
\newcommand{\supp}{\operatorname{supp}}
\DeclareMathOperator*{\argmin}{arg\,min}
\DeclareMathOperator*{\argmax}{arg\,max}
\newcommand{\entrymeta}[3]{{\sffamily\small\color{muted}\textbf{\texttt{#1}}\quad|\quad #2\par #3\par}}
\newcommand{\Problemref}[1]{\texttt{#1}}
\newcommand{\JELref}[2]{\texttt{#2} (JEL \texttt{#1})}
\begin{document}
\section*{Wages versus job amenities}
\textbf{Problem ID:} \texttt{J31-440}\quad \textbf{JEL:} \texttt{J31}\par
% BEGIN ECONOMICS STATEMENT
\label{problem:OP-0355}
\label{audited:ECON-073}
{\sffamily\small\color{muted}Original subject group: Labor markets and work\par}
\label{definitions:problem:30007637}
\textbf{Audit classification:} Background; proposed extension. Full review checked. It explains identification and confounding; its conclusion does not explicitly pose this wage-equivalent inequality problem as open.
\subsection*{Definitions and precise research target}
\textbf{Complete editorial problem.} Fix a worker population and an observable set of jobs. Job \(j\) supplies wage \(w_j\) and amenity vector \(a_j\), including scheduling, security, location, and workplace conditions. Worker \(i\) has utility \(u_i(w,a)\), strictly increasing in wage, and a feasible consideration set \(C_i\). For a declared common amenity bundle \(a_0\), define wage-equivalent job value \(e_{ij}\) by \(u_i(e_{ij},a_0)=u_i(w_j,a_j)\), when this inverse exists. Let \(j(i)\) denote actual employment. The target is the distribution of \(e_{i,j(i)}\), together with its inequality relative to the distribution of observed wages. Identify or bound the required preferences using credible offer variation, stated choices validated against behavior, and information about \(C_i\). Separate compensation for amenities from firm productivity, rent sharing, and endogenous matching. Report how conclusions depend on the reference bundle and heterogeneous preferences; interpersonal utility comparisons need additional declared normalization. A policy counterfactual equalizing amenities must additionally model wage and job reallocation. The proposed measurement task is to quantify these adjustments reliably in a specified population, not to infer their sign from competitive hedonic regressions alone.

\textbf{Definitions and admissible scope.} Let the wage domain be an interval \(\mathcal W\subseteq\mathbb R\) and amenity space \(\mathcal A\). Assume \(u_i(\cdot,a_0)\) is continuous and strictly increasing. Require \(u_i(w_j,a_j)\) to lie in its range; this guarantees a unique wage equivalent \(e_{ij}\), which need not be positive. If \(E|e_{i,j(i)}|<\infty\), define its distribution under a fixed worker measure. Use a variance comparison only under finite second moments; use Gini or top-share comparisons only when their denominators and domain conditions hold. Job choices additionally depend on consideration sets, offer probabilities, match-specific utility, and tie-breaking. Interpersonal comparison of money equivalents is conditional on the same currency, price base, and reference amenity bundle, not a utility normalization inferred from a hedonic regression. The baseline distributional target is restricted to a fixed employed cohort, with positive wages and \(j(i)\) defined for each included worker. If nonemployed persons are included, specify a separate outside-option bundle before computing its money equivalent. Choose the inequality functional as an input with a stated domain; for example, \(\operatorname{Var}(e)\) requires finite second moments. The data input includes offered and accepted jobs and experimentally or observationally identified choice probabilities. The admissible preference/consideration-set family and measurement restrictions determine which value distributions are compatible with those probabilities. Return that identified set rather than claiming that wage data alone recover \(e\).
\subsection*{Variants and assumption changes}
\begin{enumerate}
\item Experimental offers (editorial specialization informed by the review): randomize wages and one amenity over common offer support. Identify willingness to accept for that amenity under stable preferences and known offer sets; report only bounds outside the experimental wage range.
\item Equilibrium equalization (editorial): set \(a_j=a_0\) for all employers and recompute wages, costs, participation, and sorting in a declared model. Compare inequality after adjustment with the baseline money-equivalent distribution; the two are different objects.
\end{enumerate}
\subsection*{References and source audit}
\textbf{Support and access.} Full review checked. It explains identification and confounding; its conclusion does not explicitly pose this wage-equivalent inequality problem as open. \textbf{Locator:} Sections 8.2 and 9, printed pp. 60--64 and 75--76.
\begin{enumerate}\raggedright
\item\hypertarget{definitions:source:30007637:1}{} Alexandre Mas. 2025. \emph{Non-Wage Amenities}. Section 8.2, printed p. 63; conclusion, printed pp. 75–76.
\url{https://www.rfberlin.com/wp-content/uploads/2025/08/25051.pdf}.
DOI: \url{https://doi.org/10.3386/w33643}.
\end{enumerate}

\subsection*{Integrated refinement: ECON-073}
{\sffamily\small\color{muted}Valuing security and other nonwage attributes of informal jobs\par}
This refinement adopts the additional source conventions
(page~\pageref{audited:conventions}), subject to its local restrictions.
\textbf{Informal-job security and constrained offer sets.}
Specialize amenities to $(w,h,r,b,a)$: wage, hours, separation risk, benefits and
autonomy. For a security improvement $r_0\to r_1$, define compensation by
$U_i(w-CV_i,h,r_1,b,a)=U_i(w,h,r_0,b,a)$ when a unique solution exists.
Identify the distribution of this and other attribute valuations using
specified preference, offer and measurement models. Fix currency, horizon,
reference bundle and baseline worker population; specify utility monotonicity,
the income range and outside options. Accepted informal/formal jobs can reflect
constrained or selected offers rather than unconstrained preferences. Separate
wage changes, benefits and security; validated offer variation is needed to
attribute a valuation to one attribute.
\subsection*{Supplied source audit for ECON-073}
Local [S1], [S2], etc. in this audit and the references immediately below
belong to ECON-073, independently of the original entry's reference numbering.
[S1] and [S2] support valuation of nonwage work conditions. The compensation equation is editorial and is conditional on the stated bundle and offer constraints.
\subsection*{References for integrated refinement ECON-073}
{\footnotesize\raggedright
\par\noindent\textbf{[S1]} Gabriel Ulyssea, Matteo Bobba, Lucie Gadenne, and Mariaflavia Harari (2025). \emph{Informality}. VoxDevLit 6(2), April 2025.\par \textit{Verified locator / source role:} Primary publisher page checked for review identity, authors, version and background. The specific published agenda is corroborated in [S2]; the exact formal target is editorial.\par \url{https://voxdev.org/voxdevlit/informality}\par\smallskip
\par\noindent\textbf{[S2]} Oliver Hanney / VoxDev (living compilation). \emph{Evidence gaps: Key unanswered questions in development economics}. Accessed 8 October 2026. The page currently covers 20 topics; the ``16-topics URL is a historical slug.\par \textit{Verified locator / source role:} Topic heading: Informality. Supports the broad research agenda; this is not a verbatim quotation or an attribution of the displayed mathematical target.\par \url{https://voxdev.org/topic/evidence-gaps-key-unanswered-questions-across-16-topics-development-economics}\par\smallskip
}

\subsection*{Common conventions: Source mathematical conventions}

These conventions apply to each entry unless explicitly replaced.

\textbf{Models and identification.} A primitive model \(m\) specifies all probability laws, feasible choices, information, equilibrium and measurement rules used by its target. The admissible class \(\mathcal M\) is fixed independently of outcomes. If \(P_O(m)\) is its observable probability law and \(\tau(m)\) the target, its identified set at observable law \(P\) is
\[
 \mathcal I_\tau(P)=\{\tau(m):m\in\mathcal M,\ P_O(m)=P\}.
\]
Point identification means that this set is a singleton. An empty set signals incompatibility of data and assumptions. Coordinate bounds need not describe the joint set. Bounds are sharp only if every retained target is attainable or arbitrarily approachable by compatible models. Finite bounds require explicit restrictions on unobserved quantities; unrestricted confounding, hidden wealth, extrapolation or arbitrary equilibrium selection can yield uninformative sets.

\textbf{Causal interventions.} A potential outcome \(Y_i(p)\) is the outcome under a complete policy regime \(p\), including timing, financing, information and market spillovers. The target population and units are fixed before treatment. With interference, use the complete assignment vector or a justified exposure mapping; individual no-interference notation alone cannot identify an economy-wide intervention. A difference of mean potential outcomes does not require the cross-world joint distribution. Conditioning on post-treatment migration, survival, enrollment or employment changes the estimand and needs separate justification.

\textbf{Statistical statements.} An estimator, confidence set, forecast or test specifies a sampling law, model class and loss/coverage criterion. Uniform \((1-\alpha)\)-coverage means \(\inf_{m\in\mathcal M}\Pr_m\{\tau(m)\in C_n\}\geq1-\alpha\), or an explicitly asymptotic version. The whole parameter space trivially has coverage, so informativeness requires a stated width, risk or power objective. A forecasting improvement is relative to a named benchmark, information set and evaluation population.

\textbf{Equilibrium and optimization.} An equilibrium comprises feasible plans, optimality, beliefs where needed, budget constraints and all market/resource clearing. The primitives or a stated benchmark define each correspondence \(\mathcal E(p;m)\). Nonemptiness is assumed or proved, never inferred just from compact policy sets. A selected-equilibrium target uses a declared selection rule; a robust target quantifies over all permitted equilibria. Use suprema or infima unless attainment is established. An \(\varepsilon\)-optimal policy is feasible and within \(\varepsilon\) of the finite optimum value. Report an empty feasible set separately.

\textbf{Welfare and accounting.} Normative utility, interpersonal normalization, social weights, numeraire, discounting and the use of public receipts are primitives. Behavioral utility need not coincide with the welfare criterion. Household welfare includes ownership income and rents once; adding profits again to compensating variation generally double-counts them. A monetary equivalent is defined by an explicit transfer and price convention, with monotonicity and range conditions. Average effects, mechanism contributions, transfers and welfare losses are different targets. Infinite sums require convergence and dynamic feasible plans require terminal or no-Ponzi conditions.

\textbf{Source versions.} A working-paper date, revision date and journal publication year are distinguished. A DOI for a journal version is not automatically the DOI of an earlier manuscript. Locators refer to the stated version and printed pages where specified. A metadata-only check does not verify a claim about a full-text conclusion. Every entry's source subsection narrows the provenance claim accordingly.

\subsection*{Common conventions: Additional-source status and mathematical conventions}

\label{audited:conventions}

Of the 100 supplied ECON entries, 65 distinct problems appear here; 32 close
matches refine earlier entries, and three duplicate questions map to existing
entries. Appendix~\ref{app:audited-crosswalk} lists every destination.
The attachment's P/R/S/U distinctions and source audits are retained. Its
precise mathematical formalizations are editorial unless the individual audit
says otherwise. Candidate subsidy targets ECON-004--006 retain their U flags;
the resolved approval-core question ECON-007 updates OP-0202 above.
The following supplied conventions govern both this part and the attributed
ECON refinements elsewhere in the catalogue, subject to local restrictions.

\textbf{Parameterized questions.} A problem family lists inputs that must be fixed before an instance is evaluated: population, horizon, policies, information, data law, model class and welfare weights. These are required choices, not quantities the citations are assumed to specify. Undefined applications should not be treated as identified estimands.

\textbf{Indices and integrability.} Sets of agents, firms and regions are finite unless a population measure is explicitly used. \(i,j\) index units, \(r\) regions, \(t\) dates and \(h\) horizons. \(\mathbb E\) and \(\Pr\) refer to the declared population and law; \(\mathbf1\{A\}\) is an indicator. Every displayed expectation and discounted sum needs existence in the stated domain. Infinite horizons use a discount in \((0,1)\) and a convergence condition; finite horizons may use discount one.

\textbf{Optimization and existence.} A displayed \(\max\), \(\min\), or \(\arg\max\) is conditional on attainment. For finite feasible menus this is immediate; general spaces need continuity and compactness/coercivity or another existence argument. Without attainment, the target is the supremum/infimum and arbitrarily near-optimal feasible choices. \(\inf\varnothing=+\infty\). Empty feasibility and an unidentified target are different issues.

\textbf{Potential outcomes.} \(Y_i(z)\) is the outcome under a specified intervention, not the observed conditional mean among people choosing \(z\). In binary comparisons \(\Delta Y_i=Y_i(1)-Y_i(0)\). Specify assignment, treatment versions, compliance, information and observation horizon. Interference is allowed under a full policy vector; compressed exposure notation needs a justified mapping. No interference, consistency and overlap are substantive assumptions, not automatic consequences of notation.

\textbf{Populations and observation.} Fix baseline eligibility and population weights. Conditioning on post-policy employment, survival, relocation or program uptake can change the population being compared. Death is not ordinary loss to follow-up; outcomes truncated by death need a composite convention, joint survival target, or explicitly defined survivor stratum. Fertility-changing interventions need a population functional rather than an unexplained fixed descendant cohort. Measurement and missingness models must be supplied.

\textbf{Identification.} A structural model \(m\in\mathcal M\) implies observable law \(P_m\) and target \(\theta(m)\). Identification means
\[
 P_{m_1}=P_{m_2}\ \Longrightarrow\ \theta(m_1)=\theta(m_2)
 \quad\text{for all }m_1,m_2\in\mathcal M .
\]
Without identification, the target is
\[
 \Theta(P;\mathcal M)=\{\theta(m):m\in\mathcal M,\ P_m=P\}.
\]
Writing \(\theta(P)\) before establishing identification can hide incompatible latent models. Confidence sets additionally need a sampling law, confidence level, asymptotic sequence and uniformity class. Point identification, coverage and informative precision are separate properties.

\textbf{Mediation.} Total effects of randomized assignment are distinct from cross-world natural indirect effects. Belief/effort-channel effects require a meaningful mediator intervention and additional measurement, exclusion or mediation assumptions. Randomization of the initial policy alone does not supply those assumptions. Report bounds or interventional alternatives when the stronger target is unsupported.

\textbf{Equilibrium.} A counterfactual \(W(z)\), \(p(z)\), or \(\mu_t^z\) requires individual optimization, feasibility, government/resource budgets, market clearing and state-transition laws. State equilibrium existence and selection, or report ranges over compatible equilibria. Initial-state integration includes subsequent endogenous transitions; current-state integration must use the policy-dependent distribution. Reduced-form invariance after policy is not assumed.

\textbf{Welfare accounts.} Scores, utility, health, dollars and ecological quantities require explicit conversion before addition. Fees, taxes, subsidies, wages and extortion payments are transfers between parties; distributional weights can make them welfare relevant, but their face value is not automatically a resource loss. State ownership, geographic boundaries, foreign-income weights and fiscal finance. Do not add profits to household welfare already including dividends, or count land capitalization and the underlying benefit twice. A benefit-cost expression must distinguish gross benefits from net welfare and subtract every real cost exactly once.

\textbf{Derivatives and risk.} Log derivatives require positive arguments; local responses require admissible perturbations and specified equilibrium branches. Differentiation under expectations needs regularity/domination, and boundaries may allow only one-sided derivatives. Risk uses a specified law; ambiguity uses a declared nonempty law set on a common history space. Adaptive policies must be nonanticipative. A robust criterion is an editorial choice unless attributed explicitly.

\textbf{Scope overlaps.} Allocation, collective choice, contracts, household bargaining and coordinated policy overlap economics with optimization and game theory. They are retained because they appeared in the original catalogue; no claim of a strictly game-theory-free selection is made.
% END ECONOMICS STATEMENT
\end{document}
